% !TeX root = ../教材题库_全部条目_转写工作稿.tex
% Source: XJ-B1-C05-TH-P180-Y0494
% 第5章 / 5.3 三角函数的图象与性质 / 教材 p.180 / PDF 第187页
% OCR未校对：公式与文字均可能有误；未完成数学TeX转写；禁止直接发给学生
% 原题及相邻原解截图（核对用）：../题目原图/XJ-B1-C05-TH-P180-Y0494-01.webp
% 本文件只在转写完成后作为题干使用；OCR 条目尚不是可靠题干。
\begin{bankitem}{p.180 思考（含紧邻说明）}{第5章 / 5.3 三角函数的图象与性质 / 教材 p.180 / PDF 第187页}{XJ-B1-C05-TH-P180-Y0494}{OCR未校对：公式与文字均可能有误；未完成数学TeX转写；禁止直接发给学生}
\begin{ocrdraft}
[教材p.180 / PDF第187页]
\par
？思考前面认识了正弦函数y一sin工及余弦函数3，一cos的图象，你认为它
\par
们分别具有哪些性质？
\par
1.周期性
\par
由诱导公式sin(x十2差丌）=sinx(k\ensuremath{\text{[OCR:EUR]}}Z)与cos(x十2差丌）=cos工）可知，当自
\par
变量增加或减少2丌的整数倍时，
\par
与cos的值均会重复出现．为了定量地描
\par
述这种“周而复始”型的变化规律，我们引人周期函数的概念
\par
一般地，对于函数3-丆0），如果存在非零常数T，使得当工取定义域内每一
\par
个值时，x\ensuremath{\pm}T都有定义，并且
\par
则称函数3-不为周期函数，T称为这个函数的一个周期．
\par
如果T是函数3-丆）的周期，则由不到一丆0十T）一不0十T）十T)一丆陆十
\par
2T）知道2T也是它的周期．同理可知，T的所有非零整数倍都是3-丆）的周期．
\par
按照这个概念，3，=sinx,y=cos工都是周期函数，2丌及2丌的所有非零整数倍
\par
也都是它们的周期．但从图象上可以看出，比2丌更小的正数不可能是，一
\par
y=cos的周期．也就是说，这两个函数的图象向左（右）平移比2丌更短的距离不可
\par
能与原来的曲线重合，因此称2丌是y=sin3，=cos工的最小正周期．最小正周期
\par
常简称为周期（D．
\par
2\ensuremath{\cdot}值域与最值
\par
从三角函数的定义和图象可知：正弦函数、余弦函数的值域都是匚一1，0，最
\par
大值都是1，最小值都是一l.
\par
对于正弦函数3，一sin产当且仅当x一一十2丌（ez）时取得最大值1，当且仅当
\par
一十2差丌（差ez）时取得最小值一1.
\par
2
\par
(D在本书中提到三角函数的周期时，一般是指它们的最小正周期．本章中对函数最小正周期的证明均
\par
从略．
\end{ocrdraft}
\end{bankitem}
